\documentclass[11pt,oneside,reqno]{article}

\usepackage[T1]{fontenc}
\usepackage{amsmath,amssymb,amsthm}
\usepackage{mathpazo}
\usepackage[margin=1in]{geometry}
\usepackage{microtype}
\usepackage{tikz-cd}
\usetikzlibrary{arrows.meta,positioning}
\usepackage[hidelinks]{hyperref}
\hypersetup{
  pdftitle={Hochschild Homology of the b-Pseudodifferential Symbol Algebra},
  pdfauthor={Changwei Zhou}
}

\DeclareMathOperator{\HH}{HH}
\DeclareMathOperator{\gr}{gr}
\DeclareMathOperator{\End}{End}
\newcommand{\C}{\mathbb C}
\newcommand{\Z}{\mathbb Z}
\newcommand{\cA}{\mathcal A}
\newcommand{\cO}{\mathcal O}
\newcommand{\PsiL}{\Psi_{b,L}^{\mathbb Z}}
\newcommand{\Sbstar}{{}^{b}S^*X}
\newcommand{\Tbstar}{{}^{b}T^*X}
\allowdisplaybreaks

\theoremstyle{plain}
\newtheorem{theorem}{Theorem}[section]
\newtheorem{proposition}[theorem]{Proposition}
\newtheorem{lemma}[theorem]{Lemma}
\newtheorem{corollary}[theorem]{Corollary}

\theoremstyle{definition}
\newtheorem{definition}[theorem]{Definition}
\newtheorem{example}[theorem]{Example}

\theoremstyle{remark}
\newtheorem{remark}[theorem]{Remark}
\newtheorem*{discussion}{Discussion}

\title{Hochschild Homology of the\\
b-Pseudodifferential Symbol Algebra}
\author{Changwei Zhou}
\date{}

\begin{document}
\hypersetup{pageanchor=false}
\maketitle
\thispagestyle{empty}
\clearpage
\hypersetup{pageanchor=true}
\pagenumbering{roman}
\tableofcontents
\clearpage
\listoffigures
\clearpage
\pagenumbering{arabic}
\setcounter{page}{1}

\begin{abstract}
Let \(X\) be a compact \(n\)-manifold with embedded boundary and let \(E\to X\)
be a complex vector bundle. We compute the continuous Hochschild homology of
the Laurent complete-symbol algebra of classical b-pseudodifferential operators:
\[
 \cA_{b,L}(X;E)
 =
 \frac{\displaystyle\bigcup_{m\in\Z,\,\ell\geq0}
 x^{-\ell}\Psi_b^m(X;E)}
 {\displaystyle\bigcup_{\ell\geq0}
 x^{-\ell}\Psi_b^{-\infty}(X;E)}.
\]
Here \(x\) is a boundary defining function. With the natural
inductive-limit topology and completed Hochschild chains,
\[
 \HH_k^{\mathrm{cont}}(\cA_{b,L}(X;E))
 \cong H_L^{2n-k}(\Sbstar\times S^1_\sigma;\C).
\]
The circle \(S^1_\sigma\) records symbol order and \(H_L^*\) is Laurent de
Rham cohomology at the boundary. We derive the formula through the order
filtration, the Hochschild--Kostant--Rosenberg map, the Poisson differential,
and symplectic duality. We also separate this result from the Hochschild
homology of an unreduced operator algebra, which depends on the chosen
residual ideal and completion.
\end{abstract}

\section{The precise statement}

The phrase \emph{the algebra generated by b-pseudodifferential operators} is
ambiguous. At least four choices affect Hochschild homology: the allowed
operator orders, the allowed powers of a boundary defining function, the
residual ideal, and the locally convex completion. We use the standard algebra
for which the homology has a finite-dimensional geometric description.
It is the b-groupoid specialization of the Laurent symbol algebras studied in
\cite{MN,BNsymbols,BNfamilies}.

Let \(x\in C^\infty(X)\) satisfy \(x\geq0\),
\(\partial X=\{x=0\}\), and \(dx\ne0\) on \(\partial X\). Put
\[
 \PsiL(X;E)
 =\bigcup_{m\in\Z,\,\ell\geq0}x^{-\ell}\Psi_b^m(X;E),
 \qquad
 \Psi_{b,L}^{-\infty}(X;E)
 =\bigcup_{\ell\geq0}x^{-\ell}\Psi_b^{-\infty}(X;E).
\]
Both are given their natural strict inductive-limit topologies. The second
space is a two-sided ideal in the first.

\begin{definition}
The \emph{Laurent complete-symbol algebra} is
\[
 \cA_{b,L}(X;E)
 =\PsiL(X;E)/\Psi_{b,L}^{-\infty}(X;E).
\]
Its continuous Hochschild chains are
\[
 C_k^{\mathrm{cont}}(\cA_{b,L})
 =\cA_{b,L}^{\widehat\otimes(k+1)},
\]
with the completed projective tensor product and the usual Hochschild
boundary \(b\).
\end{definition}

Write
\[
 Y=\Sbstar=(\Tbstar\setminus0)/\mathbb R_+,
 \qquad Z=Y|_{\partial X}.
\]
The boundary of \(Y\) is \(Z\). The symbol-order circle is denoted
\(S^1_\sigma\).

\begin{theorem}[Hochschild homology]\label{thm:main}
For every \(k\geq0\),
\begin{equation}\label{eq:main}
 \boxed{
 \HH_k^{\mathrm{cont}}(\cA_{b,L}(X;E))
 \cong H_L^{2n-k}(Y\times S^1_\sigma;\C).}
\end{equation}
The answer is independent of \(E\). If \(X\) has one embedded boundary
hypersurface, then
\begin{align}\label{eq:expanded}
 \HH_k^{\mathrm{cont}}(\cA_{b,L}(X;E))
 \cong{}& H^{2n-k}(Y)
 \oplus H^{2n-k-1}(Y) \notag\\
 &\oplus H^{2n-k-1}(Z)
 \oplus H^{2n-k-2}(Z),
\end{align}
where a cohomology group in negative degree is zero. In particular the
homology vanishes for \(k>2n\).
\end{theorem}

The isomorphism in \eqref{eq:main} is the intrinsic statement. The four
summands in \eqref{eq:expanded} use the Laurent face decomposition and the
K\"unneth splitting for \(S^1_\sigma\).

\begin{figure}[htbp]
\centering
\begin{tikzcd}[column sep=large,row sep=large]
 \cA_{b,L}
 \arrow[r,"\text{order filtration}"]
 &\gr\cA_{b,L}
 \arrow[r,"\mathrm{HKR}"]
 &\Omega_L^*(\Tbstar\setminus0)
 \arrow[d,"\delta_{\pi_b}"]\\
 \HH_*^{\mathrm{cont}}(\cA_{b,L})
 &H_L^{2n-*}(Y\times S^1_\sigma)
 \arrow[l,"E^2=E^\infty"']
 &H_*^{\pi_b,L}(\Tbstar\setminus0)
 \arrow[l,"\text{symplectic duality}"']
\end{tikzcd}
\caption[The Hochschild-to-de Rham computation]{The computation passes from
the filtered noncommutative algebra to its commutative associated graded,
then from Poisson homology to Laurent de Rham cohomology.}
\end{figure}

\section{The b-cotangent geometry}

The Lie algebra of b-vector fields is
\[
 \mathcal V_b(X)
 =\{V\in C^\infty(X;TX):V\text{ is tangent to }\partial X\}.
\]
It is the space of smooth sections of a vector bundle \({}^bTX\to X\).
Near a boundary point, choose product coordinates
\[
 (x,y_1,\ldots,y_{n-1}).
\]
Then
\[
 \mathcal V_b(X)
 =C^\infty(X)\text{-span}
 \{x\partial_x,\partial_{y_1},\ldots,\partial_{y_{n-1}}\}.
\]
The dual local frame of \(\Tbstar\) is
\[
 \frac{dx}{x},\quad dy_1,\ldots,dy_{n-1}.
\]
Thus a b-covector has the form
\[
 \xi\frac{dx}{x}+\eta\cdot dy.
\]

On \(\Tbstar\) the tautological b-one-form and its differential are
\[
 \alpha_b=\xi\frac{dx}{x}+\eta\cdot dy,
 \qquad
 \omega_b=d\alpha_b
 =d\xi\wedge\frac{dx}{x}+d\eta\wedge dy.
\]
The form \(\omega_b\) is nondegenerate as a b-two-form. Its inverse is the
b-Poisson tensor
\[
 \pi_b=x\partial_x\wedge\partial_\xi
       +\sum_j\partial_{y_j}\wedge\partial_{\eta_j}.
\]
This is the geometric bracket appearing in the first nonzero differential of
the Hochschild spectral sequence.

\section{The associated graded algebra}

Filter \(\cA_{b,L}\) by operator order:
\[
 F_m\cA_{b,L}
 =\operatorname{image}\left(
 \bigcup_{\ell\geq0}x^{-\ell}\Psi_b^m(X;E)
 \right),
 \qquad m\in\Z.
\]
The filtration is exhaustive, separated in the complete-symbol sense, and
multiplicative. Principal symbols give
\[
 \gr_m\cA_{b,L}
 \cong \cO_L(\Tbstar\setminus0;\End E)_m,
\]
the Laurent functions that are homogeneous of degree \(m\) in the fibers.

\begin{lemma}[Morita reduction]
The coefficient bundle \(E\) does not change the Hochschild homology.
\end{lemma}

\begin{proof}
Locally, the associated graded algebra with coefficients is a matrix algebra
over the scalar symbol algebra. The local Morita maps are compatible with
changes of frame, and the standard trace and matrix-unit homotopies patch by a
partition of unity. Hence the completed Hochschild complexes for scalar
symbols and for \(\End E\)-valued symbols are quasi-isomorphic. The same
filtered argument lifts the conclusion from the associated graded algebra to
\(\cA_{b,L}(X;E)\).
\end{proof}

We therefore suppress \(E\). The Hochschild--Kostant--Rosenberg map on the
associated graded algebra is
\[
 \chi(a_0\otimes\cdots\otimes a_j)
 =\frac1{j!}a_0\,da_1\wedge\cdots\wedge da_j.
\]

\begin{proposition}[The \(E^1\)-page]
The HKR map identifies the first page of the order spectral sequence with
Laurent differential forms on \(\Tbstar\setminus0\), decomposed by fiber
homogeneity.
\end{proposition}

\begin{proof}
On each b-cotangent chart the associated graded product is ordinary
multiplication of scalar homogeneous symbols. The continuous HKR theorem
identifies its Hochschild homology with differential forms. Laurent poles in
\(x\) cause no change in the local contracting homotopy: every chain has a
uniform finite pole order, and differentiation preserves the Laurent class.
The construction is local and commutes with restriction, so a partition of
unity patches the chartwise identifications. Fiber homogeneity is preserved
throughout.
\end{proof}

\section{The first differential is Poisson homology}

For classical b-symbols \(a\) and \(c\), symbolic composition gives
\[
 \sigma_{m+m'-1}([A,C])
 =\frac1{i}\{a,c\}_b.
\]
Consequently the first differential after HKR is, up to the harmless factor
\(i\), the Poisson differential
\[
 \delta_b=[\iota_{\pi_b},d]
 =\iota_{\pi_b}d-d\iota_{\pi_b}.
\]

Let \(N=2n\) be the dimension of \(\Tbstar\). Symplectic duality is the map
\[
 *_b:\Omega_L^j(\Tbstar\setminus0)
 \longrightarrow
 \Omega_L^{N-j}(\Tbstar\setminus0)
\]
defined by contraction with the symplectic volume \(\omega_b^n/n!\).
In Darboux b-coordinates, the standard calculation gives
\begin{equation}\label{eq:symplectic-duality}
 *_b\delta_b=(-1)^{j+1}d*_b
 \quad\text{on }\Omega_L^j.
\end{equation}
It remains valid at \(x=0\), because both sides use the b-frame
\((dx/x,dy,d\xi,d\eta)\) and preserve finite Laurent pole order.

\begin{corollary}
The Poisson homology on the \(E^2\)-page is Laurent de Rham cohomology with
complementary degree:
\[
 H_j^{\pi_b,L}(\Tbstar\setminus0)
 \cong H_L^{2n-j}(\Tbstar\setminus0).
\]
\end{corollary}

\section{Homogeneity and the symbol-order circle}

Choose a positive fiber norm \(r\) on \(\Tbstar\setminus0\). Then
\[
 \Tbstar\setminus0\cong Y\times\mathbb R_+,
 \qquad (q,r)\longmapsto rq.
\]
Only total homogeneity zero survives in the de Rham cohomology of the formal
homogeneous expansion. Indeed, if \(\mathcal E=r\partial_r\) is the Euler
vector field and a form \(\alpha\) has nonzero homogeneity \(w\), Cartan's
formula gives
\[
 (d\iota_{\mathcal E}+\iota_{\mathcal E}d)\alpha
 =\mathcal L_{\mathcal E}\alpha=w\alpha.
\]
Thus \(w^{-1}\iota_{\mathcal E}\) contracts every nonzero weight.

At weight zero, a form splits uniquely as
\[
 \alpha(r,q)=\alpha_0(q)+\frac{dr}{r}\wedge\alpha_1(q).
\]
The formal order variable \(dr/r\) is represented cohomologically by the
generator \(d\theta_\sigma\) of a circle. Hence
\begin{equation}\label{eq:order-circle}
 H_L^j(\Tbstar\setminus0)_{\mathrm{hom}=0}
 \cong H_L^j(Y\times S^1_\sigma).
\end{equation}
This is the source of the extra \(S^1\) in the answer; it is not a geometric
circle in the fibers of \(\Tbstar\).

\section{Laurent cohomology and the boundary face}

Let \(M\) be a compact manifold with one embedded boundary hypersurface and
defining function \(x\). The Laurent de Rham complex is
\[
 \Omega_L^j(M)=\bigcup_{\ell\geq0}x^{-\ell}\Omega^j(M),
 \qquad d:\Omega_L^j(M)\to\Omega_L^{j+1}(M).
\]

\begin{proposition}[Face decomposition]\label{prop:face}
There is an isomorphism
\[
 H_L^j(M)
 \cong H^j(M)\oplus H^{j-1}(\partial M).
\]
The boundary summand is represented by
\(\widetilde\beta\wedge d(\chi\log x)\), where \(\beta\) is a closed
\((j-1)\)-form on \(\partial M\), \(\widetilde\beta\) is a collar extension,
and \(\chi\) is a cutoff supported in the collar and equal to one near the
boundary.
\end{proposition}

\begin{proof}
For germs at a boundary point, Taylor expansion in \(x\) reduces a Laurent
form to a finite sum of powers \(x^q\). Every term with \(q\ne0\) is exact in
the normal direction because
\[
 x^{q-1}dx=\frac1q d(x^q).
\]
The two uncontracted normal classes are \(1\) and \(dx/x=d\log x\).
Thus the cohomology sheaf of the Laurent complex has one copy of the constant
sheaf in degree zero and one boundary-supported copy in degree one. The
ordinary de Rham resolution in the tangential variables identifies their
global cohomology with \(H^j(M)\) and \(H^{j-1}(\partial M)\), respectively.
The displayed cutoff construction realizes the boundary class globally.
\end{proof}

Apply Proposition~\ref{prop:face} to \(M=Y\times S^1_\sigma\), whose boundary
is \(Z\times S^1_\sigma\). Then use
\[
 H^r(W\times S^1)
 \cong H^r(W)\oplus H^{r-1}(W).
\]
For \(r=2n-k\) this gives exactly the four summands in
\eqref{eq:expanded}.

\begin{figure}[htbp]
\centering
\begin{tikzcd}[column sep=large,row sep=large]
 H_L^r(Y\times S^1_\sigma)
 \arrow[d,equals]
 & r=2n-k\\
 H^r(Y\times S^1_\sigma)
 \oplus H^{r-1}(Z\times S^1_\sigma)
 \arrow[d,equals]\\
 \bigl(H^r(Y)\oplus H^{r-1}(Y)\bigr)
 \oplus
 \bigl(H^{r-1}(Z)\oplus H^{r-2}(Z)\bigr)
\end{tikzcd}
\caption[Interior, order, and boundary contributions]{The four summands
come from the ordinary class, the symbol-order logarithm \(d\log r\), the
boundary logarithm \(d\log x\), and their product.}
\end{figure}

For a manifold with corners, the same local argument has one generator
\(d\log x_H\) for every boundary hypersurface \(H\) through a point. If
\(\mathcal F_j(M)\) denotes the codimension-\(j\) faces, then
\begin{equation}\label{eq:corners}
 H_L^r(M)
 \cong\bigoplus_{j\geq0}
       \bigoplus_{F\in\mathcal F_j(M)}H^{r-j}(F).
\end{equation}

\section{Collapse and assembly of the computation}

Combining HKR, the Poisson differential, symplectic duality, and homogeneity
gives
\[
 E^2_k\cong H_L^{2n-k}(Y\times S^1_\sigma).
\]
The remaining input is the collapse theorem for Laurent complete symbols.

\begin{theorem}[Benameur--Nistor collapse theorem]
For the order-filtered continuous Hochschild complex of a Laurent
complete-symbol algebra associated with a b-groupoid, the spectral sequence
converges and degenerates at \(E^2\).
\end{theorem}

This is the point at which we use the convergence and degeneration theorem of
\cite{BNfamilies}; the preceding identifications make its \(E^2\)-page
explicit in the present b-setting.

This theorem is proved by first working over product charts, where the
filtered algebra is a completed formal symbol algebra, and then using the
locality of continuous Hochschild homology to glue the calculation. The
Laurent filtration supplies uniform bounds on pole orders, and the order
filtration supplies completeness in the negative direction. These are the
two points needed for convergence; without them the formal \(E^2\)-calculation
need not determine the homology of the original algebra.

The theorem identifies \(E^2\) with the associated graded of Hochschild
homology. A quantization map splits the order filtration as a topological
vector space, giving \eqref{eq:main}. Different quantizations change the
splitting but not the intrinsic \(H_L^{2n-k}\) description. This completes
the computation.

\section{Checks and consequences}

\subsection{The closed-manifold case}

If \(\partial X=\varnothing\), Laurent cohomology is ordinary cohomology and
\(Z=\varnothing\). Formula \eqref{eq:expanded} reduces to
\[
 \HH_k^{\mathrm{cont}}(\cA(X;E))
 \cong H^{2n-k}(S^*X)
 \oplus H^{2n-k-1}(S^*X),
\]
equivalently \(H^{2n-k}(S^*X\times S^1_\sigma)\). This is the
Brylinski--Getzler computation \cite{BG}.

\subsection{Degree zero and residue traces}

Since \(\HH_0(\cA)=\cA/[\cA,\cA]\), continuous traces factor through
\(\HH_0\). Theorem~\ref{thm:main} gives
\[
 \HH_0^{\mathrm{cont}}(\cA_{b,L})
 \cong H_L^{2n}(Y\times S^1_\sigma).
\]
If \(n\geq2\), \(X\) is connected, and \(\partial X\) is connected, then
\(Y\) has nonempty boundary and \(Z\to\partial X\) has connected sphere
fiber. Dimension considerations reduce the right side to
\(H^{2n-2}(Z)\cong\C\). Thus the Laurent complete-symbol algebra has one
residue trace up to scale. With several boundary components, one obtains a
boundary-face residue for each connected component of \(Z\).

\subsection{A product cylinder}

Let \(X=[0,1]\times N\), where \(N\) is closed of dimension \(n-1\), and put
\[
 W=S(\mathbb R\oplus T^*N).
\]
Then \(Y\simeq[0,1]\times W\) and \(Z=W\sqcup W\). Therefore
\[
 \HH_k^{\mathrm{cont}}(\cA_{b,L})
 \cong H^{2n-k}(W)
 \oplus H^{2n-k-1}(W)^{\oplus3}
 \oplus H^{2n-k-2}(W)^{\oplus2}.
\]
This example separates the order-circle contribution from the two boundary
faces.

\section{What the theorem does not compute}

The unreduced operator algebra fits into
\begin{equation}\label{eq:extension}
 0\longrightarrow\Psi_{b,L}^{-\infty}(X;E)
 \longrightarrow\PsiL(X;E)
 \longrightarrow\cA_{b,L}(X;E)
 \longrightarrow0.
\end{equation}
The answer for \(\PsiL\) is not obtained by simply reusing
\eqref{eq:main}. It depends on the topology placed on the residual ideal and
on whether kernels are required to vanish rapidly at the side faces, merely
have polyhomogeneous expansions, or are completed in an operator norm.

For an H-unital residual ideal, excision gives the long exact sequence
\[
 \cdots\to
 \HH_k(\Psi_{b,L}^{-\infty})
 \to\HH_k(\PsiL)
 \to\HH_k(\cA_{b,L})
 \xrightarrow{\partial}
 \HH_{k-1}(\Psi_{b,L}^{-\infty})
 \to\cdots .
\]
Melrose and Nistor compute the analogous two-filtration sequence in the cusp
calculus and explain how the same mechanism applies to the b-calculus
\cite{MN}. The
connecting maps are the algebraic location of the index and eta corrections.
Thus \eqref{eq:main} is the complete answer for Laurent complete symbols, while
\eqref{eq:extension} is the correct starting point for any specified unreduced
operator algebra.

\section{Relation to index theory}

An elliptic b-operator determines an invertible matrix over the complete-symbol
algebra. Its \(K_1\)-class pairs with cyclic cocycles obtained from the
Hochschild classes above. The ordinary symbol residue detects the interior
term; the boundary logarithmic class detects the indicial or eta term. In the
operator extension, the boundary map sends the symbol class to the Fredholm
index whenever the normal operator is invertible.

The work of Liu and Ma on differential \(K\)-theory and eta-invariant
localization supplies a geometric refinement of this pairing: differential
forms and eta data refine the topological \(K\)-class. It does not replace the
Hochschild spectral-sequence calculation \cite{LMlocalization,LMlambda}. The
logical division is therefore
\[
 \text{Melrose--Nistor and Benameur--Nistor: algebraic homology,}
\]
\[
 \text{Liu--Ma: differential-\(K\) and eta localization of index data.}
\]

\phantomsection
\addcontentsline{toc}{section}{References}
\begin{thebibliography}{99}

\bibitem{BNfamilies}
M.-T. Benameur and V. Nistor,
\emph{Homology of algebras of families of pseudodifferential operators},
J. Funct. Anal. \textbf{205} (2003), 1--36;
arXiv:math/0201198.

\bibitem{BNsymbols}
M.-T. Benameur and V. Nistor,
\emph{Homology of complete symbols and noncommutative geometry},
in \emph{Quantization of Singular Symplectic Quotients},
Progr. Math. 198, Birkh\"auser, 2001, 21--46.

\bibitem{BG}
J.-L. Brylinski and E. Getzler,
\emph{The homology of algebras of pseudodifferential symbols and the
noncommutative residue},
K-Theory \textbf{1} (1987), 385--403.

\bibitem{LMlocalization}
B. Liu and X. Ma,
\emph{Differential K-theory and localization formula for eta-invariants},
Invent. Math. \textbf{222} (2020), 545--613;
arXiv:1808.03428.

\bibitem{LMlambda}
B. Liu and X. Ma,
\emph{Lambda-ring structure in differential K-theory},
arXiv:2602.03450, 2026.

\bibitem{MelroseAPS}
R. B. Melrose,
\emph{The Atiyah--Patodi--Singer Index Theorem},
Research Notes in Mathematics 4, A K Peters, 1993.

\bibitem{MN}
R. B. Melrose and V. Nistor,
\emph{Homology of pseudodifferential operators I: manifolds with boundary},
arXiv:funct-an/9606005.

\bibitem{MoroianuNistor}
S. Moroianu and V. Nistor,
\emph{Index and homology of pseudodifferential operators on manifolds with
boundary}, in \emph{Perspectives in Operator Algebras and Mathematical
Physics}, 2008.

\end{thebibliography}

\end{document}
